\documentclass[10pt,a4paper]{amsart}
\usepackage{amsmath,amssymb,amsthm,graphicx}
% Retained minimal preamble for the Stacks Project chapter more-algebra.tex (More on Algebra).
% The chapter file pinned-original-source.tex carries only \input{preamble} on its first line, so no
% preamble text lives inside the chapter itself.  The Stacks Project uses one single shared preamble for
% every chapter; the snapshot stages it identically as sources/stacks-chow/preamble.tex and
% sources/stacks-etale-cohomology/preamble.tex.  The block below is copied mechanically, line for line,
% from sources/stacks-chow/preamble.tex lines 14-18 (the Xy-pic setup) and 154-217 (the native theorem
% declarations together with the mathematical macros).  Nothing was added, removed or redefined; the
% wrapper class and the amsmath/amssymb/amsthm/graphicx base packages are supplied by the assembler.
% Chapter: more-algebra.tex -- More on Algebra.
% For commutative diagrams we use Xy-pic
\usepackage[all]{xy}

% We use 2cell for 2-commutative diagrams.
\xyoption{2cell}

% Theorem environments.
%
\theoremstyle{plain}
\newtheorem{theorem}[subsection]{Theorem}
\newtheorem{proposition}[subsection]{Proposition}
\newtheorem{lemma}[subsection]{Lemma}

\theoremstyle{definition}
\newtheorem{definition}[subsection]{Definition}
\newtheorem{example}[subsection]{Example}
\newtheorem{exercise}[subsection]{Exercise}
\newtheorem{situation}[subsection]{Situation}

\theoremstyle{remark}
\newtheorem{remark}[subsection]{Remark}
\newtheorem{remarks}[subsection]{Remarks}

\numberwithin{equation}{subsection}

% Macros
%
\def\lim{\mathop{\mathrm{lim}}\nolimits}
\def\colim{\mathop{\mathrm{colim}}\nolimits}
\def\Spec{\mathop{\mathrm{Spec}}}
\def\Hom{\mathop{\mathrm{Hom}}\nolimits}
\def\Ext{\mathop{\mathrm{Ext}}\nolimits}
\def\SheafHom{\mathop{\mathcal{H}\!\mathit{om}}\nolimits}
\def\SheafExt{\mathop{\mathcal{E}\!\mathit{xt}}\nolimits}
\def\Sch{\mathit{Sch}}
\def\Mor{\mathop{\mathrm{Mor}}\nolimits}
\def\Ob{\mathop{\mathrm{Ob}}\nolimits}
\def\Sh{\mathop{\mathit{Sh}}\nolimits}
\def\NL{\mathop{N\!L}\nolimits}
\def\CH{\mathop{\mathrm{CH}}\nolimits}
\def\proetale{{pro\text{-}\acute{e}tale}}
\def\etale{{\acute{e}tale}}
\def\QCoh{\mathit{QCoh}}
\def\Ker{\mathop{\mathrm{Ker}}}
\def\Im{\mathop{\mathrm{Im}}}
\def\Coker{\mathop{\mathrm{Coker}}}
\def\Coim{\mathop{\mathrm{Coim}}}

% Boxtimes
%
\DeclareMathSymbol{\boxtimes}{\mathbin}{AMSa}{"02}

%
% Macros for moduli stacks/spaces
%
\def\QCohstack{\mathcal{QC}\!\mathit{oh}}
\def\Cohstack{\mathcal{C}\!\mathit{oh}}
\def\Spacesstack{\mathcal{S}\!\mathit{paces}}
\def\Quotfunctor{\mathrm{Quot}}
\def\Hilbfunctor{\mathrm{Hilb}}
\def\Curvesstack{\mathcal{C}\!\mathit{urves}}
\def\Polarizedstack{\mathcal{P}\!\mathit{olarized}}
\def\Complexesstack{\mathcal{C}\!\mathit{omplexes}}
% \Pic is the operator that assigns to X its picard group, usage \Pic(X)
% \Picardstack_{X/B} denotes the Picard stack of X over B
% \Picardfunctor_{X/B} denotes the Picard functor of X over B
\def\Pic{\mathop{\mathrm{Pic}}\nolimits}
\def\Picardstack{\mathcal{P}\!\mathit{ic}}
\def\Picardfunctor{\mathrm{Pic}}
\def\Deformationcategory{\mathcal{D}\!\mathit{ef}}

% --- appended 2026-10-12 (Chief): the project own definitions of the reference/slogan/history
% display environments, copied mechanically from the shared Stacks preamble lines 8-12
% (\usepackage{verbatim} plus the three \newenvironment lines, which the chapter preamble keeps
% verbatim from the project preamble). Nothing added, removed or redefined: these environments
% carry no mathematical content and the project itself typesets them as comments. Without these
% three lines a counted statement whose own text uses them fails with Environment slogan undefined.
\usepackage{verbatim}
\newenvironment{reference}{\comment}{\endcomment}
\newenvironment{slogan}{\comment}{\endcomment}
\newenvironment{history}{\comment}{\endcomment}
\usepackage{cleveref}
\crefname{originallemma}{original lemma}{original lemmas}
\Crefname{originallemma}{Original lemma}{Original lemmas}
\makeatletter
\newlabel{more-algebra-lemma-flattening-complete-local-universal-property}{{0528}{source}{}{sourcepointer-1}{}}
\newlabel{more-algebra-lemma-flattening-complete-local-universal-property@cref}{{[originallemma][0][]0528}{[1][0][]source}{}{}{}}
\makeatother
\crefname{originallemma}{original lemma}{original lemmas}
\Crefname{originallemma}{Original lemma}{Original lemmas}
\makeatletter
\newlabel{lemma-flattening-complete-local-noetherian}{{0526}{source}{}{sourcepointer-2}{}}
\newlabel{lemma-flattening-complete-local-noetherian@cref}{{[originallemma][0][]0526}{[1][0][]source}{}{}{}}
\makeatother
\crefname{originallemma}{original lemma}{original lemmas}
\Crefname{originallemma}{Original lemma}{Original lemmas}
\makeatletter
\newlabel{equation-flat-at-primes-over}{{052W}{source}{}{sourcepointer-3}{}}
\newlabel{equation-flat-at-primes-over@cref}{{[originallemma][0][]052W}{[1][0][]source}{}{}{}}
\makeatother
\crefname{originallemma}{original lemma}{original lemmas}
\Crefname{originallemma}{Original lemma}{Original lemmas}
\makeatletter
\newlabel{lemma-base-change-flat-at-primes-over}{{052X}{source}{}{sourcepointer-4}{}}
\newlabel{lemma-base-change-flat-at-primes-over@cref}{{[originallemma][0][]052X}{[1][0][]source}{}{}{}}
\makeatother
\crefname{originallemma}{original lemma}{original lemmas}
\Crefname{originallemma}{Original lemma}{Original lemmas}
\makeatletter
\newlabel{lemma-limit-preserving-flat-at-primes-over}{{05LM}{source}{}{sourcepointer-5}{}}
\newlabel{lemma-limit-preserving-flat-at-primes-over@cref}{{[originallemma][0][]05LM}{[1][0][]source}{}{}{}}
\makeatother
\crefname{originallemma}{original lemma}{original lemmas}
\Crefname{originallemma}{Original lemma}{Original lemmas}
\makeatletter
\newlabel{lemma-flattening-complete-local-noetherian-property-by-finite-type}{{0527}{source}{}{sourcepointer-6}{}}
\newlabel{lemma-flattening-complete-local-noetherian-property-by-finite-type@cref}{{[originallemma][0][]0527}{[1][0][]source}{}{}{}}
\makeatother
\title{A universal property of the flattening ideal for finite type algebras}
\author{The Stacks Project Authors}
\date{stacks-project:more-algebra.tex:local-snapshot-20261012; GNU FDL 1.2 or later}
\begin{document}
\maketitle
\setcounter{subsection}{1}
\noindent\textit{Editorial note.} All delivered mathematical text is the author's own, copied from the pinned Stacks Project chapter file more-algebra.tex; no statement or proof step was rewritten, reordered or removed.\par\medskip
\noindent\textit{Editorial note.} The chapter file pinned-original-source.tex carries only \textbackslash{}input\{preamble\} on its first line, so no preamble text lives inside the chapter. The retained file sources/stacks-more-algebra/retained-preamble.tex is a mechanical, line-for-line copy of the project's single shared preamble as staged in the snapshot (Xy-pic setup plus theorem declarations and mathematical macros). Nothing was added, removed or redefined, and the excerpt serves only as preamble, never as a span coordinate.\par\medskip
\noindent\textit{Editorial note.} Both counted spans start at the \textbackslash{}begin line of their environment, so the statement and the author's proof are delivered whole; surrounding prose of the section is deliberately outside the counted spans and is carried as standard prerequisites instead.\par\medskip
\noindent\textit{Editorial note.} The citation fields (source\_id, source\_version, authors, source\_url, license.name, source\_title) are identical to those of every other card this lane has produced from this chapter, so the batch selector treats them as one source.\par\medskip
\noindent\textit{Allowed standard prerequisites (editorial).}\begin{itemize}
\item Base change behaviour of the flat-at-primes-over condition. Lemma \textbackslash{}ref\{lemma-base-change-flat-at-primes-over\}, cited at proof line 4524. Gives one implication.
\item The flat-at-primes condition can be descended along a directed colimit of rings. Lemma \textbackslash{}ref\{lemma-limit-preserving-flat-at-primes-over\}, cited at proof lines 4533-4535. Moves the condition to a single essentially finite type subalgebra.
\item The finite type version of the characterisation of the flattening ideal. Lemma \textbackslash{}ref\{lemma-flattening-complete-local-noetherian-property-by-finite-type\}, cited at proof lines 4536-4538. Concludes that the ideal dies in the subalgebra.
\end{itemize}\medskip
\ifdefined\hypertarget\hypertarget{sourcepointer-1}{}\fi
\noindent\textit{Source reference.} Original lemma 0528: focused-seeds/AGP/sources/stacks-snapshot-20261011/tags.txt, source lines 6561--6561. The lemma's own canonical tag, quoted so that readers can resolve the label inside the delivered text; it points at the counted statement itself and carries no proof content.\par\medskip
\ifdefined\hypertarget\hypertarget{sourcepointer-2}{}\fi
\noindent\textit{Source reference.} Original lemma 0526: focused-seeds/AGP/sources/stacks-snapshot-20261011/tags.txt, source lines 6559--6559. Navigation pointer to a statement the counted proof refers to but does not contain; nothing load-bearing is delegated to it.\par\medskip
\ifdefined\hypertarget\hypertarget{sourcepointer-3}{}\fi
\noindent\textit{Source reference.} Original lemma 052W: focused-seeds/AGP/sources/stacks-snapshot-20261011/tags.txt, source lines 6584--6584. Navigation pointer to a statement the counted proof refers to but does not contain; nothing load-bearing is delegated to it.\par\medskip
\ifdefined\hypertarget\hypertarget{sourcepointer-4}{}\fi
\noindent\textit{Source reference.} Original lemma 052X: focused-seeds/AGP/sources/stacks-snapshot-20261011/tags.txt, source lines 6585--6585. Navigation pointer to a statement the counted proof refers to but does not contain; nothing load-bearing is delegated to it.\par\medskip
\ifdefined\hypertarget\hypertarget{sourcepointer-5}{}\fi
\noindent\textit{Source reference.} Original lemma 05LM: focused-seeds/AGP/sources/stacks-snapshot-20261011/tags.txt, source lines 7241--7241. Navigation pointer to a statement the counted proof refers to but does not contain; nothing load-bearing is delegated to it.\par\medskip
\ifdefined\hypertarget\hypertarget{sourcepointer-6}{}\fi
\noindent\textit{Source reference.} Original lemma 0527: focused-seeds/AGP/sources/stacks-snapshot-20261011/tags.txt, source lines 6560--6560. Navigation pointer to a statement the counted proof refers to but does not contain; nothing load-bearing is delegated to it.\par\medskip
\begin{lemma}
\label{lemma-flattening-complete-local-universal-property}
With notation $R \to S$, $M$, and $I$ and assumptions as in
Lemma \ref{lemma-flattening-complete-local-noetherian}.
In addition assume that $R \to S$ is of finite type.
Then for any local homomorphism of local rings
$\varphi : (R, \mathfrak m) \to (R', \mathfrak m')$
the following are equivalent
\begin{enumerate}
\item condition (\ref{equation-flat-at-primes-over}) holds
for $(R' \to S \otimes_R R', \mathfrak m', M \otimes_R R')$, and
\item $\varphi(I) = 0$.
\end{enumerate}
\end{lemma}

\begin{proof}
The implication (2) $\Rightarrow$ (1) follows from
Lemma \ref{lemma-base-change-flat-at-primes-over}.
Let $\varphi : R \to R'$ be as in the lemma satisfying (1).
As $R$ is Noetherian we see that $R \to S$ is of finite presentation
and $M$ is an $S$-module of finite presentation.
Write $R' = \colim_\lambda R_\lambda$
as a directed colimit of local $R$-subalgebras $R_\lambda \subset R'$,
with maximal ideals $\mathfrak m_\lambda = R_\lambda \cap \mathfrak m'$
such that each $R_\lambda$ is essentially of finite type over $R$. By
Lemma \ref{lemma-limit-preserving-flat-at-primes-over}
we see that condition (\ref{equation-flat-at-primes-over}) holds for
$(R_\lambda \to S \otimes_R R_\lambda, \mathfrak m_\lambda,
M \otimes_R R_\lambda)$ for some $\lambda$. Hence
Lemma \ref{lemma-flattening-complete-local-noetherian-property-by-finite-type}
applies to the ring map $R \to R_\lambda$ and we see that
$I$ maps to zero in $R_\lambda$, a fortiori it maps to zero in $R'$.
\end{proof}

\end{document}
